% 2.1 =============== 序列问题到递归问题 ===============

\section{动态规划：由序列问题转向递归形式}  \label{sec-sp-fe}

对于第一章中所要求解的动态优化问题，我们的目标是寻找一个最优值序列，通常我们把这样的问题称为\textbf{序列问题}。
但这一问题形式的求解存在一定不足：处理无限期界过于复杂，将面临无穷的差分方程序列，我们希望降维。
Stocky、Lucas等人基于贝尔曼建立的动态规划理论，将序列问题转化为递归形式求解，由此建立了现代宏观的核心方法，
而递归也成为绝大多数宏观模型的核心特征。

具体看：当计划者在第0期选择了$c_0,k_1$后，其余选择可以等到下一期做出。更一般的来说，
无穷序列问题可以分解为一期一期的小问题：对于$t+1$期，可以在第$t$期进行选择而不用在第0期进行，即定义从当期到下期的映射，
这一问题形式被称为\textbf{递归形式}（Recursive Formulation）。
\textbf{序列问题求解后，我们得到了无穷个数值（Number）；对于递归问题，我们求解所得的是一个方程（Function）,其本质是一个泛函问题}。
利用泛函方程（Function Equation）对动态优化问题进行研究就是\textbf{动态规划}（Dynamic Programmming），
这一概念最早由Richard E. Bellman提出。在本节，我们将建立一般化形式的递归问题，进而对动态规划中的基本概念进行初步介绍。


